\section{Síntesis y ampliación}

\subsection{Conclusiones centrales}

\begin{importantbox}{Doce conclusiones transferibles}
\begin{enumerate}
  \item Entre las hyperscaler primitives y la application intent sigue haciendo falta una capa de compilación.
  \item Una abstracción debe crear nuevos invariantes, no limitarse a cambiarle la apariencia a la consola.
  \item La cardinalidad y el ciclo de vida de los objetos del producto determinan si el mapeo de recursos es económico.
  \item FDI genera un IR a partir de la semántica del framework y luego lo mapea a routing, cache y functions.
  \item ISR consume origin compute según la frecuencia con que cambia el contenido, no según la frecuencia de acceso.
  \item En Build vs. Buy, el desarrollo propio debe ubicarse en la capa diferenciadora, donde el insight se acumula como interés compuesto.
  \item El immutable deployment junto con un pointer mutable hace que preview, promote y rollback sean componibles.
  \item El consumption pricing solo se convierte en una señal de retroalimentación para los desarrolladores cuando es atribuible.
  \item La telemetry debe colocar versión, rendimiento, errores y spend en una misma línea de tiempo.
  \item El soft cap protege el aprendizaje y la continuidad; el hard cap protege contra el peor gasto posible, pero puede interrumpir el negocio.
  \item Un AI workload debe distinguir entre active CPU e I/O wait; la compute density no puede depender de un número fijo de pods.
  \item El Day 1 consigue la adopción, el Day 100 absorbe la complejidad y el Day 1000 gana la confianza operativa.
\end{enumerate}
\end{importantbox}

\subsection{Tarea práctica: diseñar una framework-aware AI deployment platform}

Elige una AI web application, por ejemplo un chat multimodelo, un coding agent, un análisis de documentos o una página de comercio electrónico generativa, y completa el siguiente diseño:

\begin{enumerate}
  \item Define la semántica de aplicación que puede aportar el framework y las restricciones de negocio que la plataforma no puede inferir automáticamente;
  \item Diseña el Build Output / IR y enumera assets, functions, routes, cache, tools y policies;
  \item Presenta el modelo de immutable deployment, domain pointer, preview y rollback;
  \item Clasifica las solicitudes en cuatro tipos (CPU-bound, I/O-bound, streaming y background) y diseña la capacity policy;
  \item Diseña el telemetry schema que relacione deployment version, latency, error, model token, CPU, memory y spend;
  \item Diseña los soft/hard controls y explica umbrales, retraso de ejecución, continuidad del negocio y abuse handling;
  \item Escribe por separado las pruebas de aceptación del Day 1, el Day 100 y el Day 1000.
\end{enumerate}

\begin{knowledgebox}{Criterios de aceptación}
Una propuesta de alta calidad no se limita a decir «usar serverless con escalado automático». Debe explicar cómo entra la framework intent en el IR, cómo distingue el planificador entre active compute y wait, cómo se atribuye el costo a cada operation, cómo evita el hard control afectar por error a usos legítimos y cómo se maneja el rollback cuando se encuentra con una stateful dependency.
\end{knowledgebox}

\subsection{Lecturas adicionales}

\begin{enumerate}
  \item Vercel, \href{https://vercel.com/blog/framework-defined-infrastructure}{Framework-Defined Infrastructure}.
  \item Vercel, \href{https://vercel.com/blog/build-output-api}{Announcing the Build Output API} y \href{https://vercel.com/docs/build-output-api/v3}{Build Output API v3}.
  \item Vercel, \href{https://vercel.com/blog/life-of-a-request-application-aware-routing}{Life of a Vercel request: Application-aware routing}.
  \item Next.js, \href{https://nextjs.org/docs/app/guides/incremental-static-regeneration}{Incremental Static Regeneration}.
  \item Vercel, \href{https://vercel.com/blog/introducing-fluid-compute}{Introducing Fluid Compute} y \href{https://vercel.com/docs/fluid-compute}{current documentation}.
  \item Vercel, \href{https://vercel.com/docs/spend-management}{Spend Management} y \href{https://vercel.com/docs/observability}{Observability}.
  \item v0, \href{https://v0.app/docs}{Documentation}.
\end{enumerate}

\subsection{Ruta de lectura sugerida}

Primero lee FDI y Build Output API, y dibuja la frontera de compilación «código fuente—semántica del framework—IR—despliegue»; después lee request routing e ISR, y sigue cómo una solicitud llega a metadata, cache, function y origin. A continuación lee Spend Management y Observability, y relaciona los elementos de monitoreo con la consumption feedback vista en clase. Por último lee Fluid Compute y usa la proporción active CPU / wall time para explicar por qué las AI requests cambiaron el objetivo de planificación de serverless. Al leer cada página oficial, registra de forma uniforme \textbf{el contrato de entrada, los recursos que se generan o administran, las señales observables, la ruta de falla/rollback y las diferencias entre la versión actual y la clase histórica}; si la página no explica con claridad un failure mode, agrega por tu cuenta los posibles riesgos del plano de control, del plano de datos y de la medición del consumo.

\end{document}
